\begin{lemma}[Defective-chain structure, part (a)]
Let $(T,e_0,\varphi)$ be globally maximum, let $\alpha$ be defective, and
let the free color $\gamma$ be missing at $u\in V(T)$.  The full
$(\alpha,\gamma)$-component through $u$ is a path from $u$ to a vertex
$z\notin V(T)$.
\end{lemma}
\begin{proof}
Put $X=V(T)$ and let
$I=\noderef{TwoColorIncidenceGraph}(G,\varphi;\alpha,\gamma)$.  The
statement-level conclusion of \noderef{ScheideGlobalMaximumOrder} says
that $X$ is both \noderef{ColoredVertexSetElementary} and
\noderef{ColoredVertexSetClosed}.  Since $\alpha$ is defective, the last
clause of \noderef{ScheideGlobalMaximumBoundaryFacts} gives at least three
members of
$\noderef{BoundaryColorEdges}(G,\varphi,X,\alpha)$.  Choose one such edge
identity $b$.  By \noderef{BoundaryColorEdges},
$b\in\noderef{BoundaryEdges}(G,X)$ and $\varphi(b)=\alpha$.  Applying
closedness to $b$ at an arbitrary $x\in X$ shows that
$\alpha\notin\noderef{MissingColors}(G,\varphi,x)$.  Since
\noderef{MissingColors} is the complement of
\noderef{PresentColors}, the color $\alpha$ is present at every vertex of
$X$.

In particular, $\alpha$ is present at $u$, whereas the hypothesis says
that $\gamma$ is missing there.  Hence $\alpha\ne\gamma$: equality would
put the same color in both \noderef{PresentColors} and its complement
\noderef{MissingColors} at $u$.  We next compute the degrees in $I$.
The definition of \noderef{PresentColors} supplies an edge identity
incident with a vertex for each color present there, where incidence is
\noderef{Incident}.  There is at most one incident edge identity of any
fixed color.  Indeed, two distinct such identities would be
\noderef{AdjacentEdges}, while the \noderef{GlobalMaximumTriple}
hypothesis contains a \noderef{KTriple}; its
\noderef{IsProperEdgeColoring} clause forbids equal colors on two such
edges.

At $u$ there is consequently exactly one incident $\alpha$-edge and no
incident $\gamma$-edge, the latter assertion being exactly what it means
for $\gamma$ to belong to \noderef{MissingColors} at $u$.  By the
adjacency definition in \noderef{TwoColorIncidenceGraph}, these are all
neighbors of the vertex-node $u$, so $u$ has degree one in $I$.  Now fix
$x\in X\setminus\{u\}$.  Elementariness says that the missing-color sets
at $u$ and $x$ are disjoint.  Since $\gamma$ is missing at $u$, it is not
missing at $x$, and hence is present there.  We already proved that
$\alpha$ is present at $x$.  The preceding uniqueness argument gives one
incident edge of each color; because $\alpha\ne\gamma$, these two edge
identities are distinct.  They are precisely the neighbors of $x$ in
$I$.  Thus every vertex-node in $X\setminus\{u\}$ has degree two in $I$.

Let
$C=\noderef{TwoColorComponent}(G,\varphi;\alpha,\gamma,u)$.  Every
edge-node $f$ in $C$ also has degree two in $I$.  To see this explicitly,
$u$ has a neighbor by the preceding paragraph, so $C$ is not a singleton.
Because $f$ is connected to $u$, it cannot be isolated; an adjacency at
$f$ and the definition of \noderef{TwoColorIncidenceGraph} show that
$\varphi(f)$ is $\alpha$ or $\gamma$.  Write the unordered endpoint pair
of $f$ as $\{r,s\}$.  The loopless axiom in
\noderef{LooplessMultigraph} gives $r\ne s$.  By
\noderef{Incident}, both endpoint vertex-nodes are adjacent to $f$ in
$I$, and the definition of $I$ says that there are no other neighbors of
an edge-node.  Hence $f$ has exactly the two distinct neighbors $r$ and
$s$.  This also agrees with the universal upper bound
\noderef{TwoColorIncidenceDegreeAtMostTwo}.

Properness having been obtained above, apply
\noderef{TwoColorComponentSpanningPath} to $u$.  It gives a simple path
$Q$ in $I$ whose support is exactly $C$.  The node $u$ belongs to $C$ and
therefore lies on $Q$.  It cannot be internal on $Q$, since an internal
node of a simple path has two distinct path-neighbors, whereas $u$ has
degree one.  Thus $u$ is an end of $Q$.  Moreover, $Q$ has positive
length: the unique $\alpha$-edge at $u$ is a second node of $C$, and $Q$
spans $C$.  Reverse $Q$ if necessary and call the resulting path $P$, so
that $P$ starts at $u$; let $w\ne u$ be its other end.

We claim that $w$ has degree one in $I$.  The last edge of the
positive-length path gives one neighbor $r$ of $w$.  Suppose that $w$ had
a second neighbor $y\ne r$.  Since adjacency to $w\in C$ puts $y$ in the
same connected component, the support equality for $P$ puts $y$ on $P$.
The node $y$ is not $w$, since $I$ is a simple graph and hence has no
loop at $w$.  It also cannot be $u$: if it were, then $w$ would be a second
neighbor of $u$ in addition to the first edge of $P$ (and when $P$ has
length one, $r=u$, contrary to $y\ne r$), contradicting the already proved
degree-one assertion at $u$.  Thus $y$ is internal on $P$.  Its two
distinct path-neighbors are different from $w$, because the only path-node
adjacent along $P$ to the end $w$ is its predecessor $r$, while
$y\ne r$.  Together with the extra neighbor $w$, the node $y$ would have
at least three neighbors.  This contradicts
\noderef{TwoColorIncidenceDegreeAtMostTwo}.  Therefore $r$ is the only
neighbor of $w$, proving the claim.

The end $w$ is not an edge-node, since every edge-node of $C$ has degree
two.  Hence $w$ is a vertex-node, say $w=z$.  It is not a vertex of $X$:
it is not $u$, and every other vertex-node of $X$ has degree two.  We have
therefore obtained a simple incidence path $P$ from $u$ to a vertex
$z\notin X$ whose support is exactly $C$.  Together with the hypothesis
$u\in X$, these are precisely the four clauses of
\noderef{ScheideDefectiveChainPath}, so $z$ and $P$ are the required
witnesses.
\end{proof}
